% !TEX root = ../presentation.tex
% =============================================================
%  Frame — the 500 kV reference line, worked numbers
%  Part of ac-transmission; \input from ac-transmission.tex,
%  after reactance.tex and before power-transmission.tex.
%
%  Collects the worked values that the preceding frames refer to, so
%  those frames carry only the physics.  Same line as
%  ohl-configuration.tex and series-compensation.tex.  All values from
%  fig/angular_stability_figures.py (exact two-port) and the nominal pi
%  of two-port.tex.  P_max of each model occurs at delta = theta_B.
% =============================================================

\begin{frame}{The 500\,kV reference line}
  \begin{block}{Worked example}
    {The same line carries the whole chapter\\
    The model you choose moves $P_{max}$ by up to 20\,\%.}
  \end{block}

  \vspace{0.5em}
  \begin{columns}[T,onlytextwidth]
    \begin{column}{0.42\textwidth}
      \renewcommand{\arraystretch}{1.05}
      {\footnotesize
      \begin{tabular}{@{}l l r@{}}
        \textbf{Line} & & \multicolumn{1}{c}{\textbf{Value}} \\
        \hline
        $U_n$        & kV            &                500 \\
        $f$          & Hz            &                 60 \\
        $l$          & mi (km)       &          300 (483) \\
        Conductor    &               & $3\times$ Bluebird \\
        Geometry     & ft / in       &    35 flat / 18 \\
        \hline
        $R'$         & m$\Omega$/km  &               56.9 \\
        $X'$         & $\Omega$/km   &              0.336 \\
        $C'$         & nF/km         &               12.7 \\
        \hline
        $Z_w$        & $\Omega$      &                266 \\
        $P_{nat}$    & MW            &                939 \\
        \hline
      \end{tabular}}
    \end{column}
    \begin{column}{0.55\textwidth}
      \renewcommand{\arraystretch}{1.05}
      {\footnotesize
      \begin{tabular}{@{}l l r@{}}
        \textbf{At 300\,mi} & & \multicolumn{1}{c}{\textbf{Value}} \\
        \hline
        $\underline{Z}$     & $\Omega$ & $27.5+j\,162.2$ \\
        $\underline{Y}$     & mS       & $j\,2.31$ \\
        Charging at 1.0\,p.u. & MVAr   & 578 \\
        $\beta\cdot l$      & $^\circ$ & 35.2 \\
        \hline
      \end{tabular}}\\[0.8em]
      {\footnotesize
      \begin{tabular}{@{}l r r r r@{}}
        \textbf{Model} & \multicolumn{1}{c}{\textbf{$\left|A\right|$}}
          & \multicolumn{1}{c}{\textbf{$\vartheta_B$}}
          & \multicolumn{1}{c}{\textbf{$\left|B\right|$}}
          & \multicolumn{1}{c}{\textbf{$P_{max}$}} \\
         & & \multicolumn{1}{c}{$^\circ$}
          & \multicolumn{1}{c}{$\Omega$} & \multicolumn{1}{c}{MW} \\
        \hline
        Short, lossless & 1.000 & 90.0 & 162.2 & 1541 \\
        Nominal $\pi$   & 0.813 & 80.4 & 164.5 & 1266 \\
        \atrowhighlight Distributed & 0.819 & 81.0 & 154.4 & 1364 \\
        \hline
      \end{tabular}}
    \end{column}
  \end{columns}

  \centering
  {\color{atGray}\footnotesize $P_{max}$ occurs at $\delta=\vartheta_B$;
  the simplification moves it to $90^\circ$ and overstates it by 13\,\%.}
\end{frame}
